% Bagian: first-order-logic
% Bab: introduction
% Subbagian: soundness-completeness

\documentclass[../../../include/open-logic-section]{subfiles}

\begin{document}

\olfileid[id]{fol}{int}{scp}

\olsection{Kesahihan dan Kelengkapan}

Kita juga akan memperkenalkan sistem !!{derivation} bagi logika orde pertama.
Ada banyak sistem !!{derivation} yang telah dikembangkan oleh para ahli
logika, tetapi semuanya mendefinisikan relasi !!{derivability} yang sama di
antara !!{sentence}s. Kita mengatakan bahwa $\Gamma$ \emph{!!{derive}s}~$!A$,
$\Gamma \Proves !A$, jika terdapat !!a{derivation} dari jenis tertentu yang
didefinisikan secara tepat. !!^{derivation}s selalu merupakan susunan simbol
yang berhingga---mungkin berupa daftar !!{sentence}s, atau suatu struktur yang
lebih rumit. Tujuan sistem !!{derivation} ialah menyediakan alat untuk
menentukan apakah suatu !!a{sentence} diakibatkan oleh suatu himpunan~$\Gamma$.
Untuk mencapai tujuan tersebut, harus berlaku bahwa $\Gamma
\Entails !A$ jika, dan hanya jika, $\Gamma \Proves !A$.

Jika $\Gamma \Proves !A$ tetapi tidak berlaku bahwa $\Gamma \Entails !A$,
sistem !!{derivation} kita terlalu kuat: sistem itu membuktikan terlalu
banyak. Sifat bahwa jika $\Gamma \Proves !A$ maka $\Gamma \Entails !A$
disebut \emph{kesahihan}, dan merupakan persyaratan minimal bagi setiap sistem
!!{derivation} yang baik. Sebaliknya, jika $\Gamma \Entails !A$ tetapi tidak
berlaku bahwa $\Gamma \Proves !A$, sistem !!{derivation} kita terlalu lemah:
sistem itu tidak membuktikan cukup banyak. Sifat bahwa jika $\Gamma \Entails
!A$ maka $\Gamma \Proves !A$ disebut \emph{kelengkapan}. Kesahihan biasanya
relatif mudah dibuktikan (melalui induksi pada struktur !!{derivation}s, yang
didefinisikan secara induktif). Kelengkapan lebih sulit dibuktikan.

Kesahihan dan kelengkapan mempunyai sejumlah konsekuensi penting. Jika suatu
himpunan !!{sentence}s~$\Gamma$ !!{derive}s suatu kontradiksi (seperti $!A
\land \lnot !A$), himpunan itu disebut \emph{tak konsisten}. $\Gamma$ yang
tak konsisten tidak dapat mempunyai model apa pun; himpunan itu tidak dapat
dipenuhi. Dari kelengkapan diperoleh arah sebaliknya: setiap $\Gamma$ yang
bukan merupakan himpunan tak konsisten---atau, sebagaimana akan kita katakan,
\emph{konsisten}---
mempunyai suatu model. Bahkan, pernyataan ini ekuivalen dengan kelengkapan,
dan merupakan bentuk kelengkapan yang benar-benar akan kita buktikan. Ini
merupakan hasil yang mendalam dan mungkin mengejutkan: fakta bahwa Anda tidak
dapat membuktikan $!A \land \lnot !A$ dari $\Gamma$ menjamin adanya
!!a{structure} yang sesuai dengan deskripsi~$\Gamma$. Jadi, kelengkapan
memberikan jawaban bagi pertanyaan: himpunan !!{sentence}s mana yang mempunyai
model? Jawabannya: semua dan hanya himpunan yang konsisten.

Teorema kesahihan dan kelengkapan mempunyai dua konsekuensi penting: teorema
kekompakan dan teorema L\"owenheim--Skolem. Keduanya merupakan hasil penting
dalam teori model, dan dapat digunakan untuk menetapkan banyak hasil yang
menarik. Kita sudah menyebutkan dua di antaranya: logika orde pertama tidak
dapat menyatakan bahwa !!{domain} suatu !!a{structure} bersifat hingga atau
bahwa domain tersebut !!{nonenumerable}.

Secara historis, semua hal ini---cara mendefinisikan sintaksis dan semantik
logika orde pertama, cara mendefinisikan sistem !!{derivation} yang baik, cara
membuktikan bahwa sistem tersebut sahih dan lengkap, serta memahami dengan
jelas apa yang dapat dan tidak dapat diungkapkan dalam bahasa orde pertama---
memerlukan waktu yang lama untuk ditemukan dan dirumuskan dengan benar. Kini
kita mengetahui cara melakukannya, tetapi menelusuri semua perinciannya masih
dapat membingungkan dan menjemukan. Namun, hal itu juga penting, karena metode
yang dikembangkan di sini bagi bahasa formal logika orde pertama diterapkan
di berbagai bidang dalam logika, ilmu komputer, dan linguistik. Jadi,
mempelajari perinciannya akan bermanfaat dalam jangka panjang.

\end{document}
